\section{Endpoint determination and meromorphic continuation}
\label{sec:continuation}

The analytic parametrization constructed in
\Cref{thm:joint-parametrization} will now be continued in its phase
variable.  The main conclusion of this section is an alternative of a
particularly rigid kind: at a finite first horizontal wall of
meromorphy, there is no continuation patch at even one point of the
wall.  The proof combines two ingredients.  A finite billiard chain is
algebraically determined by derivatives of its endpoints, and a local
branch change cannot be compatible with the even glancing
parametrization.  Both endpoint formulations will be needed: one with
lines as endpoints and one with impact points as endpoints.
Complex extensions of billiard action-angle data also enter the
perturbative rigidity argument of \cite{KS2018}; the first-wall
alternative and endpoint determination below are proved here.

\subsection{Joint geometry and horizontal widths}

Forget the orientation of the line with coordinates
\((\Theta(y,t),p(y,t))\), and denote the resulting line in the complex
projective plane by \(L(y,t)\).  Define
\begin{equation}
 X(y,t)=L(y-t/2,t)\cap L(y+t/2,t).
 \label{cont:impact-definition}
\end{equation}
For real, sufficiently small \(t\ne0\), this is a boundary impact.  The
reflection equations identify it with the support point whose normal
angle is the mean of the two incident line angles.  Consequently the
apparent singularity in \eqref{cont:impact-definition} at \(t=0\) is
removable, and
\[
 X(y,0)=\Gamma(y),
\]
where \(\Gamma\) is the boundary parametrized by the normalized
coordinate of the preceding sections.  Both \(L\) and \(X\) are even in
\(t\), are \(2\pi\)-periodic in \(y\), and are real-symmetric.

We treat projective maps as meromorphic data: in any fixed generic
affine chart their coordinates are meromorphic functions.  An
intersection or a reflection can therefore be used wherever its
defining rational expressions are not identically degenerate.  The
identities
\begin{equation}
 X(y,t)\in L(y-t/2,t)\cap L(y+t/2,t),\qquad
 \rank \dd X\le1
 \label{cont:incidence-rank}
\end{equation}
and the reflection identity at the tangent \(X_y\) hold initially and
continue as differential-algebraic identities.  Here and below the
complex Euclidean product is bilinear, not Hermitian.  In particular,
reflection in a nonisotropic tangent is rational in that tangent and
the incident line.  All required nondegeneracies hold near the real
collar and thus hold generically on any continuation under
consideration.

Choose a basic strip \(\{|\Im y|<s_*\}\) strictly narrower than the
initial common domain.  Then choose a parameter disk \(B\) centered at
zero sufficiently small compared with \(s_*\).  These choices are kept
fixed in this section.

\begin{definition}
\label{cont:width-definition}
For \(F=L\) or \(F=X\), let \(h_F(t)\) be the supremum of the heights
\(b>s_*\) for which the primary branch of \(F\) has a single-valued
meromorphic extension to
\[
 \{-s_*<\Im y<b\}\times V
\]
for some open neighborhood \(V\subset B\) of \(t\), agreeing with the
initial branch on the basic strip.  The neighborhood \(V\) may depend
on \(b\).  The value \(+\infty\) is allowed.
\end{definition}

The functions \(h_F\) are lower semicontinuous: an extension to height
\(b\) works on an open parameter neighborhood.  Uniqueness on the basic
strip makes the extensions agree.  Compactness in \(y\) always means
compactness modulo the period \(2\pi\).

\begin{lemma}
\label{cont:width-relation}
After the parameter disk has been reduced if necessary,
\begin{equation}
 h_L(t)=h_X(t)+\frac{|\Im t|}{2}.
 \label{cont:width-equality}
\end{equation}
\end{lemma}

\begin{proof}
Intersection in \eqref{cont:impact-definition} gives
\[
 h_X(t)\ge h_L(t)-\frac{|\Im t|}{2}.
 \]
For the opposite inequality suppose first that \(\Im t\ge0\).  The
line through \(X(y-3t/2,t)\) and \(X(y-t/2,t)\) is the preceding
billiard line.  Reflecting it at the second point gives \(L(y,t)\).
Both points, and the tangent at the second one, are available whenever
\(\Im y<h_X(t)+\Im t/2\), at each strict subheight and on a
sufficiently small parameter neighborhood.  This rational formula
agrees with the primary line in the basic strip.  The initially wider
domain supplies any lower-strip overlap needed by the shifts.
For \(\Im t<0\), use the succeeding two impacts and backward
reflection instead.  This proves the opposite inequality and hence
\eqref{cont:width-equality}.
\end{proof}

\subsection{A finite-type continuation lemma}

In the endpoint arguments, the sampled occurrences of the support
or graph function are treated as independent one-variable germs.
Holding the endpoint data fixed gives a weighted relation among
variations of these germs. The variations are initially formal;
to compare them as the impacts coalesce, we need convergence and
continuation to the coalescing edge. The following lemma supplies
both conclusions. Its final hypothesis says that all sampling
coordinates approach the same coordinate \(a\) at that edge.

\begin{lemma}[Continuation of a finite web relation]
\label{cont:web-continuation}
Let \((a,\delta)\) range over a small real rectangle
\(I\times(0,\delta_*)\).  Suppose that \(\psi_i,w_i\), \(1\le i\le m\),
and \(B\) are analytic in a neighborhood of that rectangle,
that the \(\psi_i\) are real-valued on the real rectangle,
that the \(w_i\) are nonzero, and that the differentials
\(\dd\psi_i\) have pairwise distinct directions.  Suppose
formal one-variable germs \(f_i\) at one interior basepoint
satisfy
\begin{equation}
 \sum_{i=1}^m w_i(a,\delta)f_i(\psi_i(a,\delta))=B(a,\delta).
 \label{cont:web-equation}
\end{equation}
Then these germs converge and continue along the rectangle.
If, in addition,
\[
 \psi_i(a,\delta)=a+O(\delta),\qquad
 \partial_a\psi_i(a,\delta)=1+O(\delta)
 \]
analytically up to \(\delta=0\), each continued one-variable
germ extends to that edge near every smaller interior
subinterval of \(I\).  No boundedness of the coefficients of
the induced differential system at \(\delta=0\) is required.
\end{lemma}

\begin{proof}
Locally choose a coordinate direction transverse to all the
finitely many foliations; write it as the \(a\)-direction.
Put
\[
 A_i=\partial_a\psi_i,\quad B_i=\partial_\delta\psi_i,
 \quad \rho_i=B_i/A_i,\quad
 F_{i,l}=f_i^{(l)}(\psi_i).
 \]
The \(A_i\) are nonzero and the \(\rho_i\) are distinct.
Differentiate \eqref{cont:web-equation} by
\(\partial_a^{m-r}\partial_\delta^r\), \(0\le r<m\).
The terms containing \(F_{i,m}\) have coefficient matrix
\[
 \big(w_i A_i^m\rho_i^r\big)_{\substack{0\le r<m\\1\le i\le m}},
 \]
which is invertible by the Vandermonde determinant.
All remaining terms are linear in \(F_{i,l}\), \(l<m\),
with analytic coefficients and an analytic inhomogeneous
term.  Thus each \(F_{i,m}\) is an analytic affine-linear
expression in the finite vector
\[
 V=(F_{i,l}:1\le i\le m,\ 0\le l<m).
 \]
The chain rules
\[
 \partial_a F_{i,l}=A_iF_{i,l+1},\qquad
 \partial_\delta F_{i,l}=B_iF_{i,l+1}
 \]
then give a full first-order system
\begin{equation}
 \partial_aV=M_aV+c_a,\qquad
 \partial_\delta V=M_\delta V+c_\delta.
 \label{cont:web-system}
\end{equation}

To prove convergence, solve the first coordinate linear ODE on
the base section and then the second with analytic dependence
on the first coordinate.  The Taylor series of this analytic
solution is the unique formal solution of those two selected
ODE problems.  It therefore agrees with the original formal
vector.  The residual in the unused equation of
\eqref{cont:web-system} is analytic and has zero Taylor
series, so it vanishes.  The web relation and all its
differential identities hold near the basepoint.

Linear ODEs with analytic coefficients continue along finite
coordinate segments in the interior.  Solving first along a
horizontal section and then along vertical segments, and
using analytic uniqueness on overlaps, gives continuation
throughout any smaller rectangle.  The other equations of
\eqref{cont:web-system} persist by analytic continuation.
Coordinate changes handle portions where a different
transverse direction was used.  In particular the continued
\(F_{i,0}\) still depends locally only on \(\psi_i\).

For the edge assertion, shrink the interval \(I\) and fix a
small positive section \(\delta=\delta_1\).  For \(\psi\) in
a neighborhood of the smaller interval, the equation
\(\psi_i(a,\delta)=\psi\) has a solution \(a=a_i(\psi,\delta)\)
inside the rectangle for every \(0<\delta\le\delta_1\).
Along that path,
\[
 \frac{\dd}{\dd\delta}F_{i,0}(a_i(\psi,\delta),\delta)
 =
 -\frac{B_i}{A_i}A_iF_{i,1}+B_iF_{i,1}=0.
 \]
Thus its value is represented by the analytic one-variable
germ on the fixed positive section.  That germ extends to
\(\delta=0\), regardless of any singular behavior of
\(M_a,M_\delta,c_a,c_\delta\) there.
\end{proof}

\subsection{Determination from two endpoint lines}

For a meromorphic field \(F(y,t)\), its differential field
means the field generated by its affine coordinates and all their
\(y,t\) derivatives, together with constant scalars.  Whenever
algebraicity is asserted, it is asserted coordinatewise in a
generic affine chart.

\begin{theorem}[Line-endpoint determination]
\label{thm:lens-lines}
Fix an integer \(k\ge2\), and put \(L_i=L(y+it,t)\).
Every intermediate line \(L_i\), \(0<i<k\), is algebraic over
the differential field generated by \(L_0,L_k\).
For each fixed \(k\), finitely many endpoint derivatives suffice.
The same assertion holds for translated index intervals.
All derivatives of the intermediate fields are algebraic over
the corresponding endpoint differential field.
\end{theorem}

\begin{proof}
We work first near real glancing trajectories.  The proof has
four parts: elimination of high support jets, convergence of
formal variations, isolation by a coalescing limit, and
specialization back to meromorphic functions.

\paragraph{The critical value and its endpoint data.}
Use the initial angle \(a\) and the total angular increment
\(\delta\) as local parameters.  This is a nonsingular change
from \((y,t)\): at \(t=0\), the two nonzero diagonal factors
are \(\phi'(y)\) and \(k\phi'(y)\).  Write
\[
 \alpha_0=a,\quad \alpha_k=a+\delta,\qquad
 \psi_i=\frac{\alpha_{i-1}+\alpha_i}{2},\quad
 d_i=\frac{\alpha_i-\alpha_{i-1}}2.
 \]
The intermediate angles are stationary for
\begin{equation}
 \mathcal A_{\rm line}
 =\crit_{\alpha_1,\ldots,\alpha_{k-1}}
 \sum_{i=1}^k2h(\psi_i)\sin d_i.
 \label{cont:line-critical-value}
\end{equation}
Its derivatives with respect to the two endpoint angles are
\(-p_0,p_k\).  Thus all its derivatives of positive order are
known from endpoint jets.  This statement uses only algebraic
extensions and algebraic operations on those jets.
Indeed, oriented unit normals and signed distances are
obtained from the unoriented projective lines by square
roots near the real branch; half-angle exponentials require
only further radicals.  Differentiating angles and changing
the two derivations then involve rational expressions in
these algebraic data and their derivatives.

On a small rectangle \(a\in I\), \(0<\delta<\delta_*\),
analyticity of the glancing family gives
\begin{equation}
 \alpha_i=a+\frac{i}{k}\delta+O(\delta^2),\qquad
 \psi_i=a+\frac{i-\tfrac12}{k}\delta+O(\delta^2).
 \label{cont:line-coalescence}
\end{equation}
Consequently the gradients of the \(\psi_i\) have distinct
directions and \(\partial_a\psi_i\ne0\).  The stationary
Hessian is invertible.  To see its sign and scale, subtract
the telescoping integral \(\int_a^{a+\delta}h\).
The leading summands become
\[
 -\frac{r(a)}{24}(\alpha_i-\alpha_{i-1})^3.
 \]
At equal steps their Hessian is negative definite on
Dirichlet variations of the interior angles.  The
differentiated analytic remainders preserve invertibility
for positive sufficiently small \(\delta\), uniformly on
a smaller a-interval.

\paragraph{Finite algebraic core.}
Fix a generic interior basepoint.  Allow the \(k\) occurrences
of the support in \eqref{cont:line-critical-value} to be
different germs \(h_i\).  The finite core consists of the
interior half-angle exponentials and the jets of every
\(h_i\) through order
\[
 m_0=\max(2,k)
 \]
at its sampled angle.  Impose base stationarity and equality
of every positive-order critical-value derivative with the
given endpoint data.  Formal critical-point Taylor
coefficients are rational functions of the support jets:
at each order one inverts the same finite stationary
Hessian and performs rational operations.

For an order \(m>m_0\), the new support jets \(h_i^{(m)}\)
enter the m-th critical-value derivative tensor with
coefficients
\begin{equation}
 2\sin d_i\,(\nabla\psi_i)^{\otimes m}.
 \label{cont:line-high-jet}
\end{equation}
Indeed a profile change beginning in order \(m\) moves the
formal stationary point only in order \(m-1\).  Its
critical-value correction beyond direct substitution has
order at least \(2m-2>m\).  The direct substitution gives
\eqref{cont:line-high-jet}.
Choose tensor components with \(0,\ldots,k-1\)
\(\delta\)-derivatives.  After dividing nonzero column
factors, their matrix is Vandermonde in
\((\partial_\delta\psi_i)/(\partial_a\psi_i)\).
It is invertible.  The gradients themselves depend only
on second-order core jets, by differentiated stationarity.
Thus every higher support jet is rationally expressible
in the core and the endpoint data.

The denominators are products of powers of a fixed finite
collection of nonzero quantities: the stationary Hessian,
the coordinate normalizations, the weights \(\sin d_i\),
the chosen coordinate derivatives, and the differences of
gradient ratios.  Work in the open locus where they are
nonzero.  After elimination, the remaining requirements are
polynomial equations in the finite core, with coefficients
in the endpoint data field.  Noetherianity reduces this
possibly infinite list to finitely many equations on that
open locus.

\paragraph{Analytic continuation of a hypothetical deformation.}
Suppose the actual core were not isolated.  Slicing a
positive-dimensional algebraic component to a curve and
parametrizing a branch gives a nonconstant complex analytic
arc of admissible cores, with parameter \(\eps\).
Every eliminated jet is holomorphic along this arc near
zero.  Move the sampled Taylor centers back to their
original locations.  Since the center changes are
\(O(\eps)\), formal translation is well-defined in the two
formal variables: each coefficient uses finitely many
shift powers.  The resulting support germs have the form
\begin{equation}
 h_i^\eps=h+\sum_{j\ge p}\eps^j f_{i,j},
 \qquad p>0,
 \label{cont:line-formal-deformation}
\end{equation}
with at least one nonzero \(f_{i,p}\).  If all support
variations vanished, formal implicit uniqueness for the
nondegenerate stationary point would also keep the angle
core fixed, contradicting nonconstancy of the arc.
Matching all positive-order derivatives of the critical
value means that its change is
\[
 \sum_{j\ge p}K_j\eps^j,
 \]
where each \(K_j\) is constant in the endpoint parameters.

We next show that every coefficient \(f_{i,j}\), initially
formal, is analytic and extends to the coalescing edge.
Induct on \(j\).  First variation at the j-th order gives
\begin{equation}
 \sum_{i=1}^k2\sin d_i\,f_{i,j}(\psi_i)
       =K_j-B_j(a,\delta).
 \label{cont:line-variation-web}
\end{equation}
The term \(B_j\) depends only on lower perturbation orders
and on the original stationary branch.  By the induction
hypothesis and Hessian inversion it is analytic throughout
the positive rectangle.  Apply
\Cref{cont:web-continuation} with the weights \(2\sin d_i\)
and the foliations \(\psi_i\).  It proves convergence,
continuation, and extension to \(\delta=0\) near a smaller
interior a-interval.  Notice that this is coefficientwise
in \(j\); no assertion of convergence of the full
\(\eps\)-series is needed.

\paragraph{The coalescing stationary limit.}
Fix an interior value of \(a\).  Put
\[
 \alpha_i=a+\delta x_i,\qquad
 z_i=x_i-x_{i-1},\qquad \sum_{i=1}^kz_i=1,
 \qquad \eps=e\delta^{2/p}.
 \]
Subtract the unperturbed integral of \(h\) and divide the
action by \(\delta^3\).  Coefficientwise in \(e\), its
stationary value tends to
\begin{equation}
 G(e^p)=\crit_{\sum z_i=1}
 \sum_{i=1}^k\left(-\frac{r(a)}{24}z_i^3
                    +e^p f_{i,p}(a)z_i\right),
 \qquad z_i(0)=1/k.
 \label{cont:line-limit-functional}
\end{equation}
Here direct Taylor expansion gives the limit of the
action and of its derivatives near the equal-step point.
A perturbation of order \(j>p\) carries the factor
\(\delta^{2j/p-2}\) and therefore disappears.  The
zeroth-order scaled Hessian tends to an invertible one.
Recursion for the formal stationary coefficients thus
gives coefficientwise convergence of the stationary
point and its value.

The coefficient of \(e^j\) contributed by the matched
critical value is \(K_j\delta^{2j/p-3}\).  The limit in
\eqref{cont:line-limit-functional} depends only on
\(x=e^p\).  At \(j=p\), boundedness forces \(K_p=0\);
at \(j=mp\), \(m\ge2\), the exponent \(2m-3\) is
positive.  Therefore every positive x-coefficient of
\(G\) vanishes, and \(G\) is constant.

Write \(f_i=f_{i,p}(a)\) and \(r=r(a)\) in the remainder
of this paragraph.  Differentiating the critical value
gives \(\sum_i f_i z_i=0\).  If \(\Lambda\) is its
Lagrange multiplier, stationarity and weighted summation
give
\[
 -\frac r8z_i^2+xf_i=\Lambda,\qquad
 \Lambda=3G=-\frac{r}{8k^2}.
 \]
Taking the roots with constant term \(z_i(0)=1/k\) and
summing them yields
\begin{equation}
 \sum_{i=1}^k
 \sqrt{1+\frac{8k^2}{r}xf_i}=k.
 \label{cont:line-power-sums}
\end{equation}
Every coefficient of positive degree in the square-root
series is nonzero.  Hence \(\sum_i f_i^m=0\) for every \(m\ge1\).
Newton's identities make all elementary symmetric
functions of the \(f_i\) vanish; each \(f_i\) is zero.
Varying \(a\) contradicts the first nonzero coefficient
in \eqref{cont:line-formal-deformation}.  The actual core
is therefore isolated.

\paragraph{From isolated values to algebraic germs.}
An isolated point of a complex algebraic set defined
over a field \(K\) has coordinates algebraic over \(K\).
For example, pass to an algebraic closure of \(K\).
Its zero-dimensional components remain zero-dimensional
after further extension of algebraically closed fields,
and an isolated point must lie in such a component.

To apply this fact to functions, do not specialize at an
arbitrary point.  Begin over a countable field containing
\(\Q(\ii)\) and the fixed numerical constants used in
the coordinates.  Adjoin the endpoint coordinate
functions, all their derivatives, the actual local
radical branches, and also the actual core functions.
The resulting function field is countable.  Choose the
real basepoint outside the poles and the zero sets of
every nonzero member of this field.  This is possible:
a nonzero analytic function cannot vanish on an open
real rectangle, and the excluded countable union has
measure zero.  Evaluation is then injective on the
field.

The algebraic equations used above are defined over the
evaluated endpoint field, with its stated algebraic
extensions.  Isolation makes the evaluated core
algebraic over that field.  Lift the coefficients of
such a polynomial relation back to functions.
Injectivity of evaluation, also after adjoining the
actual core, makes the lifted polynomial identity
hold as a germ.  The contact momentum formulas recover
the offsets of the intermediate lines from the core.
Eliminating the endpoint radical extensions proves the
asserted algebraicity over the original unoriented
endpoint differential field.  Only finitely many
endpoint derivatives occur in the resulting polynomial.
These identities persist under continuation.

Finally, in characteristic zero an algebraic element is
separable.  Differentiating its minimal polynomial
expresses each derivative algebraically over the
differential field.  This proves the assertion about
jets, and translation of \((y,t)\) proves the assertion
for other index intervals.
\end{proof}

\subsection{Determination from two endpoint impacts}

The impact formulation needs one additional scalar datum: the
area coefficient from \eqref{eq:area-coefficient}.  It supplies
the mixed derivatives of the critical length, which the two
endpoint positions alone do not directly give.

\begin{theorem}[Point-endpoint determination]
\label{thm:lens-points}
Let \(k\ge2\), and put \(P_i=X(y+it,t)\), \(0\le i\le k\).
Every intermediate point and every edge of this chain is
algebraic over the differential field generated by \(P_0,P_k\)
and the analytic coefficient \(A(t)\) in
\[
 \dd p\wedge\dd\Theta=A(t)\,\dd y\wedge\dd t.
\]
For each fixed \(k\), finitely many jets suffice.  The assertion
persists under translation and analytic continuation.
\end{theorem}

\Needspace{5\baselineskip}
\begin{proof}\leavevmode
\paragraph{The length functional and its mixed derivatives.}
Near a real glancing configuration choose orthogonal coordinates
in which the boundary is \(Y=F(x)\), with \(F''\ne0\), and
orient the arc so that \(x\) increases along the chain.  Use
the endpoint abscissas
\[
 a=x(P_0),\qquad b=x(P_k)=a+\delta
\]
as parameters.  This change of coordinates is nonsingular
near glancing.  Write \(x_i=x(P_i)\).  On a small rectangle
\(a\in I\), \(0<\delta<\delta_*\),
\begin{equation}
 x_i=a+\frac{i}{k}\delta+O(\delta^2).
 \label{cont:point-coalescence}
\end{equation}
Allow the interior occurrences of the graph to be independent
functions \(F_i\), while keeping \(F_0=F_k=F\) fixed.  Consider
the critical value
\begin{equation}
 \mathcal A_{\rm point}(a,b)=
 \crit_{x_1,\ldots,x_{k-1}}
 \sum_{i=1}^k
 \sqrt{(x_i-x_{i-1})^2+
       (F_i(x_i)-F_{i-1}(x_{i-1}))^2}.
 \label{cont:point-critical-value}
\end{equation}
The roots are the positive real lengths at the original
configuration and are continued with that choice.

The endpoint graph jets are determined by the point jets,
since \(x\) is a regular coordinate.  All mixed derivatives
of \(\mathcal A_{\rm point}\) are determined as well.  Indeed,
if \(e_{\theta_0}\) is the first unit chord velocity, endpoint
differentiation gives
\[
 (\mathcal A_{\rm point})_a
       =-P_{0,a}\cdot e_{\theta_0}.
\]
The endpoint momentum identity then gives
\begin{equation}
 (\mathcal A_{\rm point})_{ab}\,\dd a\wedge\dd b
       =\dd p_0\wedge\dd\theta_0
       =A(t)\,\dd y\wedge\dd t.
 \label{cont:point-mixed-area}
\end{equation}
The first edge is \(L(y+t/2,t)\).  Its half-shift in \(y\)
preserves \(\dd y\wedge\dd t\), so the area coefficient is
exactly the same \(A(t)\).  Equation
\eqref{cont:point-mixed-area}, the Jacobian of the endpoint
coordinate change, and differentiation recover every derivative
having at least one \(a\)- and one \(b\)-derivative.

\paragraph{Elimination of higher graph jets.}
At the original configuration subtract the telescoping graph
arc lengths between successive endpoints.  Taylor expansion of
chord length minus arc length gives the leading summands
\begin{equation}
 -\frac{F''(a)^2}{24(1+F'(a)^2)^{3/2}}
       (x_i-x_{i-1})^3.
 \label{cont:chord-arc-expansion}
\end{equation}
For completeness, if \(v=F'(a)\), \(r=F''(a)\), and
\(g=1+v^2\), the cubic coefficient of a chord over an interval
of length \(s\), after subtraction of the integral of
\(\sqrt{1+F'^2}\), is \(-r^2/(24g^{3/2})\).
This follows by expanding the square root of
\(s^2+(vs+rs^2/2+F'''(a)s^3/6)^2\) through degree three;
the \(F'''(a)\)-terms cancel against the arc integral.
At equal positive steps the Hessian of the sum in
\eqref{cont:chord-arc-expansion}, on interior variations
with fixed endpoints, is negative definite.  Thus the
stationary Hessian is invertible for small positive \(\delta\).

Take as finite core the \(x_i\), the jets of \(F_i\) through
order \(m_0=\max(2,k)\), \(0<i<k\), and the nonzero length
roots.  Impose base stationarity, the algebraic equations
for the roots, and equality of every mixed critical-value
derivative with the given endpoint data, including the
low-order derivatives.  In order \(m>m_0\), the new graph jets enter
the critical-value derivative tensor through
\begin{equation}
 w_i\,(\dd x_i)^{\otimes m},\qquad 0<i<k,
 \label{cont:point-high-jet}
\end{equation}
where \(w_i\) is the vertical component of the incoming unit
velocity minus the outgoing unit velocity.  The same
stationary-value differentiation as in
\eqref{cont:line-high-jet} proves this formula: movement of
the critical point contributes only in orders greater than
\(m\).  The weights are nonzero because the tangent of our
graph is nonvertical and the real incidence is nonzero.

Only mixed tensor components are known here, but they
suffice.  In \((a,b)\)-coordinates,
\[
 x_{i,a}=1-i/k+O(\delta),\qquad
 x_{i,b}=i/k+O(\delta).
\]
Both derivatives are nonzero and their ratios are distinct.
Use the components with \(1,\ldots,k-1\) derivatives in \(b\)
and the remaining derivatives in \(a\).  Since \(m>k\), all
are mixed components.  Their coefficient matrix is an
invertible Vandermonde matrix after nonzero column factors
have been removed.  Hence each higher graph jet is a
rational function of the finite core and the given mixed
value derivatives.  The gradients of the \(x_i\) are
obtained by inverting the second-order stationary Hessian,
so they too are already rational in this data.

Restrict to the finite open locus where the length roots,
Hessian, weights, coordinate derivatives and Vandermonde
differences are nonzero.  All denominators in the recursive
computation are products of these fixed quantities.
Elimination and Noetherianity leave an algebraic set of
possible cores defined by finitely many equations on this
open locus.  We prove that the
actual core is isolated.

\paragraph{Continuation of the variation coefficients.}
A nonisolated core would admit a nonconstant analytic arc.
Translate the moving Taylor centers as in the line proof.
The resulting formal graph variations have the form
\begin{equation}
 F_i^\eps=F+\sum_{j\ge p}\eps^j f_{i,j},
 \qquad f_{0,j}=f_{k,j}=0,
 \label{cont:point-formal-deformation}
\end{equation}
with at least one nonzero first coefficient \(f_{i,p}\).
Equality of all mixed value derivatives says precisely that
the positive-order change of critical value is
\[
 \sum_{j\ge p}\eps^j\bigl(U_j(a)+V_j(b)\bigr),
\]
initially as a formal identity.  At order \(j\), first
variation yields
\begin{equation}
 \sum_{i=1}^{k-1}w_i f_{i,j}(x_i)-U_j(a)-V_j(b)
                  =B_j(a,b),
 \label{cont:point-variation-web}
\end{equation}
where the right side is determined by lower orders.
There are \(k+1\) one-variable unknowns in this web.  The
foliations \(x_1,\ldots,x_{k-1},a,b\) have pairwise distinct
gradient directions by \eqref{cont:point-coalescence};
their weights \(w_i,-1,-1\) are nonzero.  A generic coordinate
direction is transverse to all of them, so the full
Vandermonde argument of \Cref{cont:web-continuation} applies.
Induction on \(j\) proves that all \(f_{i,j},U_j,V_j\) are
analytic on the positive rectangle and extend to the
coalescing edge near an interior \(a\)-interval.  For the
edge assertion use \((a,\delta)\): all the foliation
coordinates, including \(b=a+\delta\), are \(a+O(\delta)\)
and have \(a\)-derivative \(1+O(\delta)\).  No bound on the
web system at that edge is needed.

\paragraph{The scaled critical length.}
Fix an interior \(a\), and abbreviate
\[
 r=F''(a),\quad g=1+F'(a)^2,\quad f_i=f_{i,p}(a),\quad
 x_i=a+\delta u_i,\quad z_i=u_i-u_{i-1},\quad x=e^p.
\]
Thus \(f_0=f_k=0\), \(\sum_i z_i=1\), and \(r\ne0\).
Subtract the unperturbed graph arc from the length, multiply
by \(g^{3/2}/\delta^3\), and substitute
\(\eps=e\delta^{2/p}\).  The coefficientwise stationary
limit is
\begin{equation}
 G(x)=\crit_{\sum z_i=1}
 \sum_{i=1}^k\left[
 -\frac{r^2z_i^3}{24}
 -\frac{rx(f_i+f_{i-1})z_i}{2}
 +\frac{x^2(f_i-f_{i-1})^2}{2z_i}\right],
 \quad z_i(0)=1/k.
 \label{cont:point-limit-functional}
\end{equation}
Here are the individual terms in this limit.  The cubic
term is \eqref{cont:chord-arc-expansion}.  The constant
vertical first-variation coefficient of an edge is
\(F'(a)/\sqrt g\).  Summed against the graph offsets it
telescopes to zero because the endpoints are fixed.  The
next coefficient is
\[
 \frac{r\delta(u_i+u_{i-1})}{2g^{3/2}}.
\]
Summation by parts gives
\[
 \sum_i(u_i+u_{i-1})(f_i-f_{i-1})
       =-\sum_i(f_i+f_{i-1})z_i,
\]
which produces the displayed linear term.  The second
vertical derivative of ordinary edge length gives the
inverse-\(z_i\) quadratic term.  Higher powers of the
vertical perturbation vanish in the limit.  Terms of
perturbation order \(j>p\) have, after the constant
telescoping part has canceled, a factor
\(\delta^{2j/p-2}\) in their linear contribution and
vanish as well.  These expansions hold with derivatives
in the stationary variables near \(z_i=1/k\).  The limiting
Hessian at \(x=0\) is invertible, so formal stationary
recursion justifies the stated coefficientwise limit.

For each multiple \(j=mp\ge2p\), the separated matched
value contributes
\[
 g^{3/2}\delta^{2j/p-3}
      \bigl(U_j(a)+V_j(a+\delta)\bigr),
\]
which tends to zero.  The limit depends only on \(x=e^p\).
It follows that
\begin{equation}
 G(x)=G_0+g_1x,\qquad G_0=-\frac{r^2}{24k^2}.
 \label{cont:point-affine-limit}
\end{equation}
Let \(\lambda\) be the stationary multiplier.  Differentiating
the functional in \eqref{cont:point-limit-functional} gives
\[
 \lambda=-\frac{r^2z_i^2}{8}
       -\frac{rx(f_i+f_{i-1})}{2}
       -\frac{x^2(f_i-f_{i-1})^2}{2z_i^2}.
\]
Multiply by \(z_i\), sum, and compare with the derivative
of the critical value in \(x\).  The weighted homogeneity
of its three terms gives
\[
 \lambda=3G-2xG'=3G_0+g_1x.
\]
Consequently
\begin{align}
 \left(\frac{rz_i}{2}
       +\frac{x(f_i-f_{i-1})}{z_i}\right)^2
       &=-6G_0-2x(g_1+rf_{i-1}),\notag\\
 \left(\frac{rz_i}{2}
       -\frac{x(f_i-f_{i-1})}{z_i}\right)^2
       &=-6G_0-2x(g_1+rf_i).
 \label{cont:point-root-identities}
\end{align}
Take each square root with constant term \(r/(2k)\).
Adding both roots for each edge and then all edges gives
\(r\sum_i z_i=r\).  Equivalently,
\begin{equation}
 \sqrt{1-\beta xg_1}
 +2\sum_{i=1}^{k-1}\sqrt{1-\beta x(g_1+rf_i)}
 +\sqrt{1-\beta xg_1}=2k,
 \qquad \beta=\frac{8k^2}{r^2}.
 \label{cont:point-power-sums}
\end{equation}
Every positive power sum of the multiset
\[
 g_1+rf_0,\quad
 g_1+rf_1,g_1+rf_1,\quad\ldots,\quad
 g_1+rf_{k-1},g_1+rf_{k-1},\quad g_1+rf_k
\]
therefore vanishes.  Newton's identities make every member
zero.  The fixed endpoints give \(g_1=0\), and then
\(f_i=0\) for every \(i\).  Varying \(a\) contradicts the
first nonzero term of \eqref{cont:point-formal-deformation}.
The core is isolated.

\paragraph{Return to the endpoint function field.}
Use a countable field containing the endpoint functions,
their derivatives, the derivatives of \(A(t)\), the actual
length roots, and the actual core.  Choose the real
basepoint outside the zeros and poles of all its nonzero
members.  Evaluation is injective there.  Isolation over
the evaluated endpoint field gives polynomial algebraicity
of the core values, and injectivity lifts those polynomial
identities to germs, exactly as in the last part of
\Cref{thm:lens-lines}.  Eliminating the finitely many radical
extensions leaves the original endpoint differential
field.  The intermediate graph values give all points;
joining consecutive points gives every edge by rational
projective operations.  Differentiation of separable
polynomials gives the assertion for jets.  This completes
the point-endpoint determination.
\end{proof}

\subsection{Product continuation and comparison of widths}

Endpoint determination supplies polynomial equations, but leaves open
which branches continue from the primary collar. A hypothetical
meromorphic patch in both variables will give such equations across
a whole horizontal row for generic nearby irrational real parameters.
The monodromy argument below will turn those equations into a uniform
width improvement at those parameters. To transfer that improvement
back to the original parameter, we shall use
superharmonic comparison for \(h_L\) and \(h_X\). We prove that
comparison next, starting with the product-continuation principle
also used later for families of analytic disks.

\begin{lemma}[Product continuation]
\label{lem:meromorphic-hartogs}
Let \(D_w\) be a connected plane disk and \(0<\rho<R\).  Suppose that
\(F(w,\zeta)\) is meromorphic on
\[
 D_w\times\{\rho<|\zeta|<R\}.
 \]
If it extends meromorphically to
\(D_0\times\{|\zeta|<R\}\) for some nonempty open subdisk
\(D_0\subset D_w\), then it extends meromorphically to the whole
product \(D_w\times\{|\zeta|<R\}\).
There is also a holomorphic version: if the original data and the
initial extension are holomorphic, the resulting extension is
holomorphic.
\end{lemma}

\begin{proof}
Fix \(w_0\in D_w\).  There is an intermediate circle
\(|\zeta|=r\), \(\rho<r<R\), on a neighborhood of which
\((w-w_0)^m F\) is jointly holomorphic for some integer \(m\ge0\).
To see this, write local meromorphic quotients on a compact
subannulus at \(w=w_0\).  Remove from their denominators the powers
of \(w-w_0\) that divide them identically along that slice.  The
remaining denominator zeros on the slice are locally isolated in
\(\zeta\).  A circle avoiding these zeros exists, and a finite cover
of that circle gives one common power \(m\).

On a sufficiently small disk \(D'\) about \(w_0\), let
\(c_{-j}(w)\), \(j\ge1\), be the negative Laurent coefficients of
\((w-w_0)^mF\) on this circle.  They are holomorphic in \(w\).
Suppose that \(D'\) meets an open set of parameters already reached
by meromorphic extension to the full \(\zeta\)-disk.  For generic
parameters in that open set, the slice has finitely many poles
inside the circle.  Multiplication by a polynomial in \(\zeta\)
removes them.  Equivalently, the Hankel matrix
\[
 \big(c_{-(i+j+1)}(w)\big)_{i,j\ge0}
\]
has finite rank there.

The rank bound can be made independent of \(w\).  For each fixed
integer \(r_0\), the condition that all minors of size \(r_0+1\)
vanish is an analytic condition on \(D'\).  Unless it holds
identically, its solution set is contained in the discrete zero
set of one nonzero minor and is therefore countable.  A countable
union of such sets cannot contain the open set of generic
extended slices.  Hence all minors of one fixed size vanish
identically on \(D'\).

Over the field of meromorphic functions on \(D'\), finite Hankel
rank gives a finite recurrence for the negative coefficients.  Here
is the elementary reason.  A dependence among the first finitely
many columns with nonzero last coefficient expresses that last
column in terms of its predecessors.  Shifting the row index
gives the same dependence for all later columns.  Thus, for a
nonzero polynomial \(Q(w,\zeta)\) with meromorphic coefficients
in \(w\), every negative Laurent coefficient of
\(Q(w,\zeta)(w-w_0)^mF\) vanishes.  Clear the finitely many
coefficient denominators locally.  The remaining Laurent series
has a holomorphic extension inside the circle, given by its
Cauchy integral on that circle.  Dividing back by the polynomial
and by the cleared factors gives the desired meromorphic
extension on a neighborhood of \(w_0\) times the disk.

The set of reached parameters is therefore open and relatively
closed in \(D_w\): at a limit point the preceding local argument
uses any overlapping reached open set.  It is nonempty and hence
is all of \(D_w\).  The extensions glue by uniqueness on their
annular overlaps, and the original annulus completes them to
radius \(R\).

In the holomorphic case every negative Laurent coefficient itself
vanishes on \(D_0\).  The identity theorem makes it vanish on
\(D_w\); the same Cauchy integral gives a holomorphic extension
without a polynomial denominator.
\end{proof}

\begin{corollary}[Continuation along a family of disks]
\label{cont:disc-sweep}
Suppose a continuous path in one complex center variable has a neighborhood
on which a family of analytic disks is holomorphic in the center
and disk variables.  Assume that the pulled-back primary data are
jointly meromorphic on a common boundary collar of these disks,
that these collar germs agree on overlaps, and that a nonempty
open set of starting centers has a full-disk extension.  The
extensions then propagate along the center path.  The same
statement holds holomorphically.  In that case a uniform
boundary modulus bound, or a uniform bound on the real or
imaginary part of a scalar holomorphic component, is preserved
inside each disk.
\end{corollary}

\begin{proof}
Cover the compact part of the center path under consideration by
finitely many coordinate disks.  On each one the pulled-back
data satisfy \Cref{lem:meromorphic-hartogs}; the full disks
already reached give the nonempty open seed for the next one.
Uniqueness on the boundary collars gives agreement.  In the
holomorphic case, additional holomorphic parameters can be retained:
the same
Laurent coefficients and Cauchy integrals are holomorphic in
them, and the identity theorem applies on their common
neighborhoods.  In the holomorphic case the final assertions are
the maximum principle and its harmonic real-part version.
\end{proof}

\begin{proposition}
\label{prop:width-superharmonic}
Both \(h_L\) and \(h_X\) satisfy superharmonic comparison on \(B\).
They are positive lower semicontinuous superharmonic functions,
with the identically infinite case admitted.
\end{proposition}

\begin{proof}
Let \(t_0+r\overline{\mathbb D}\subset B\).  Choose a
nonnegative trigonometric polynomial \(q\) satisfying
\[
 q(\zeta)<h_F(t_0+r\zeta),\qquad |\zeta|=1,
 \]
and let \(g\) be a holomorphic polynomial whose boundary
imaginary part is \(q\).  Thus
\[
 \Im g(0)=\frac1{2\pi}\int_0^{2\pi}q(e^{\ii\theta})\,\dd\theta.
 \]
For real \(x_0\), \(0\le\lambda\le1\), and a small complex
variable \(w\), consider
\begin{equation}
 t=t_0+r\zeta,\qquad y=x_0+\lambda g(\zeta)+w.
 \label{cont:width-disks}
\end{equation}
On \(|\zeta|=1\), these points lie strictly below the appropriate
width.  Lower semicontinuity and compactness in
\((x_0\bmod 2\pi,\lambda,\zeta)\) provide a common annular
neighborhood in \(\zeta\) and a common disk \(|w|<d_w\) on
which the pulled-back functions are jointly meromorphic.

At \(\lambda=0\), the whole disks are in the initial basic
strip.  Increase \(\lambda\) in fixed sufficiently small
increments satisfying
\[
 \sup_{|\zeta|\le1+d_\zeta}
 |(\lambda_{\rm new}-\lambda_{\rm old})g(\zeta)|<d_w/2.
 \]
The full-disk data at the old value provide a full-disk
extension at the new value for \(|w|<d_w/2\).  Product
continuation fills the original \(w\)-disk.  Repeat finitely
many times.

For each \(\lambda\), uniqueness on the annulus glues the
extensions as \(x_0\) varies.  Adjacent values of \(\lambda\)
also glue, because their overlapping strips contain the
full-disk seeds just described.  At \(\zeta=0\) these are
strips with ordered centers
\(\Im y=\lambda\Im g(0)\).  Consequently the continued branch
is attached to the basic strip throughout the ascent and
reaches slightly above height \(\Im g(0)\) at \(t_0\).
We obtain
\[
 h_F(t_0)\ge \frac1{2\pi}
 \int_0^{2\pi}q(e^{\ii\theta})\,\dd\theta.
 \]
Approximate the positive lower semicontinuous boundary
function from below by continuous functions and then, with
a strict margin, by trigonometric polynomials.  This gives
the supermean inequality, including infinite integrals.
Together with lower semicontinuity it gives the asserted
superharmonic comparison.
\end{proof}

\subsection{No patch at a first wall}

A joint patch means a meromorphic continuation on a neighborhood
in both \(y\) and \(t\), agreeing as joint germs with the primary
branch below the wall.

\begin{theorem}[No-patch property]
\label{thm:no-patch}
Let \(F=L\) or \(F=X\), and let \(t_0\in B\cap\R\).
If \(h_F(t_0)=H_1<\infty\), no joint meromorphic patch of the
primary branch passes through any point of \(\Im y=H_1\).
The corresponding assertion at the lower wall follows by
real symmetry.
\end{theorem}

\begin{proof}
We first prove the assertion for \(L\).  Suppose a patch exists.
Choose a horizontal arc and a vertical band about \(H_1\) on
which joint data are known for all \(t\) in a neighborhood of
\(t_0\).  The patch agrees across the bottom of the band with
the lower primary tube.  Fix a narrower concentric band and
\(\eta>0\) such that \(H_1+\eta\) lies strictly inside it.
These choices are independent of the nearby irrational real
shifts used below.  Shrink the parameter neighborhood so
that a common strip below \(H_1\), reaching the bottom of
the narrower band, is available jointly.

\paragraph{Endpoint equations across the full band.}
Fix a nearby real \(t_1\) with \(t_1/(2\pi)\notin\Q\).
Forward and backward translates of any real phase meet the
good arc.  On a small phase box choose integers \(m<0<n\)
so that both \(L(y+mt,t)\) and \(L(y+nt,t)\) are supplied
by the patch throughout the narrower vertical band.
Compactness modulo \(2\pi\) gives finitely many boxes
covering the whole band.  Only finitely many shifts occur,
so these assertions hold for \(t\) in a complex neighborhood
of \(t_1\), possibly depending on \(t_1\).  Shrink this
neighborhood to retain a common bottom strip on which
every shifted field is the primary one.

Apply \Cref{thm:lens-lines} on that bottom strip in each box.
It gives nonzero polynomial equations for the coordinates
of \(L(y,t)\), with coefficients meromorphic across the
full band.  Use squarefree equations with nonzero leading
coefficients and discriminants.  These remain nonzero as
germs: their coefficients are continued actual endpoint
jets, so a nonzero identity cannot become identically zero
during continuation.  The primary branch below any first
wall in the band satisfies the equations, by vertical
continuation from the bottom strip.

\paragraph{A first obstruction has nontrivial monodromy.}
The exceptional sets are the coefficient poles and the
zeros of the leading coefficients and discriminants.
They are locally analytic curves.  Outside countably
many parameter values, a slice meets their reduced union
at isolated points where it is a smooth graph over \(t\).
Indeed, use a countable cover by analytic boxes.  Vertical
components give countably many constant parameter values;
on nonvertical components, singularities and critical
points of projection are discrete.  Intersections of
distinct components are also discrete.  Exclude all their
parameter images.  Relatively compact boxes leave only
finitely many exceptional points in a periodic row.

Choose a real irrational \(t_2\) outside those exceptions.
Suppose \(h_L(t_2)\le H_1+\eta\), so its first wall is
strictly inside the band.  Away from exceptional points,
following simple roots gives joint passage through the
wall.  Near an exceptional point the only exceptional
curve is \(y=z(t)\).  Off it the roots have finite
continuation on a punctured disk times a disk, whose
fundamental group is generated by a small \(y\)-loop
about \(z(t_2)\).  If the actual branch has trivial
monodromy on that loop, it is single-valued on the product.
After denominators are cleared, the leading coefficient
is a unit times a power of \(y-z(t)\).  A polynomial root
bound therefore gives
\[
 |L_{\rm coord}(y,t)|\le C|y-z(t)|^{-N}
\]
for a finite \(N\).  Multiplication by that power and
removal of a bounded singularity gives joint meromorphy.
If every wall point had trivial monodromy, these
extensions would glue on overlaps just below the wall.
A finite cover of the periodic wall would then continue
the tube to a strictly greater height on a parameter
neighborhood, contrary to the definition of \(h_L(t_2)\).
Some small loop must act nontrivially.

\paragraph{A defect affects only finitely many shifted indices.}
Put \(L_j=L(y+jt,t)\).  At a fixed point of the row, a
small loop can change \(L_j\) only if translation by
\(jt_2\) aligns its center with an exceptional point.
There are finitely many such points modulo period and
\(j\mapsto jt_2\pmod{2\pi}\) is injective.  Only finitely
many indices can be aligned.  Precisely, for each prescribed
finite set of indices one can choose the loop sufficiently
small to fix all its unaligned members.  No single loop
fixing an infinite collection is asserted.

First choose the loop sufficiently small for the whole finite
set of aligned indices. Further shrinking preserves its
monodromy on those fields. Among the aligned indices with
nontrivial monodromy, take the smallest and translate the
center along the orbit so that this index is \(1\). Choose
\(K\) so that every aligned index now has absolute value at
most \(K\). Choose \(b>K\), \(b\ge2\), and any further finite
collection of far indices, including both signs when needed
below. Shrink the loop to fix the unaligned members of that
collection and the unaligned members of \(\{-1,0\}\). Its
action on aligned fields is unchanged, so \(L_{-1},L_0\)
return unchanged. It produces
\begin{equation}
 M_1\ne L_1
 \label{cont:alternate-line}
\end{equation}
while fixing \(L_{-1},L_0,L_b\) and all selected far fields.

\paragraph{Descent to glancing.}
Below the wall at the real parameter \(t_2\), all selected
primary fields are single-valued jointly near every strict
subheight.  Descend there and then move along real phase
and parameter to the original collar.  A connected open
region around this connection carries all selected primary
data meromorphically, so the outward path is available.
The small loop avoids every shifted exceptional point
relevant to the finite selection.

On return \(M_1\) is algebraic over \(L_0,L_b\) and their
jets, by transporting \Cref{thm:lens-lines} around the loop.
In the connected lower region these endpoints are back
on their single-valued primary branches.  Continue the
alternative root down along paths avoiding the proper
analytic exceptional sets of its squarefree equation.
Complex detours, including an initial move off such a
set within the starting germ, are allowed.  They do not
change the selected primary data.

Near a generic point \((y_0,0)\), only \(t=0\) can be an
exceptional divisor for that polynomial.  Other analytic
curves meet this divisor discretely unless they contain it.
On a sufficiently small bidisk, a power cover \(t=\tau^d\)
kills the finite monodromy.  The polynomial root bound
makes the root meromorphic in \((y,\tau)\).  Thus \(M_1\)
has a joint meromorphic Puiseux germ at glancing.
Its inequality from \(L_1\) persists, since the primary
\(L_1\) is single-valued in the lower region.

\paragraph{The alternate germ is even.}
Put \(Z=L_0\cap M_1\).  Continued incidence, rank and
reflection identities show that \(Z\) has rank at most
one and that reflection in its tangent relates the two
lines.  It differs as a germ from the unchanged preceding
intersection \(L_{-1}\cap L_0\): continuing back to the
original real table would otherwise identify the preceding
and following intersections.

On the Puiseux cover there is a projective limit
\[
 Z^0(y)=\lim_{\tau\to0}Z(y,\tau^d),
\]
holomorphic at generic \(y\).  It is nonconstant, because
it lies on \(L(y,0)\) and the tangents of the original
strictly convex arc have no common projective point.
At a regular point of \(Z^0\), rank one implies that the
nearby image of \(Z\) lies on the same fixed analytic arc
traced by \(Z^0\).  This arc cannot lie at infinity or have
identically isotropic tangent: in either case it would be
a fixed projective line, and continuation would force
the original real table arc into that line. Indeed, the line equation
would hold on this rank-one impact branch and continue back along the
reverse path to the original impact \(L_0\cap L_1\), contradicting
strict convexity.

If \(L(y,0)\) were generically tangent to the fixed arc
at \(Z^0(y)\), differentiated incidence and tangency would
identify \(Z^0\) with the envelope of those lines,
namely \(\Gamma(y)\).  The fixed arc would be the primary
boundary germ.  Near its simple quadratic tangency,
\(L_0(y,t)\) has exactly the two primary nearby
intersections with that germ, distinct for generic small
\(t\ne0\).  The preceding intersection has been excluded;
reflection at the following one gives \(L_1\), contrary
to \eqref{cont:alternate-line}.

The fixed arc is therefore transverse to \(L(y,0)\) at a
generic point.  Intersect it with the even line \(L_0(y,t)\)
by the ordinary holomorphic implicit function theorem,
then reflect in its nonisotropic tangent.  Both operations
are holomorphic and uniquely determined near that point.
They produce \(M_1\), and prove
\begin{equation}
 M_1(y,-t)=M_1(y,t)
 \quad\hbox{as an ordinary holomorphic germ at }t=0.
 \label{cont:alternate-even}
\end{equation}

\paragraph{A continuation operation reflecting indices.}
The incoming leg can now end at a regular \((y_0,0)\)
itself; take the outgoing primary leg from the same point.
The resulting closed continuation path is denoted \(A\).
Let \(I\) denote pullback by \(t\mapsto-t\).  Perform,
in order,
\begin{equation}
 A,\qquad I,\qquad A^{-1},\qquad I.
 \label{cont:monodromy-operation}
\end{equation}
The distinguished field follows the succession
\[
 L_1\longmapsto M_1\longmapsto M_1
     \longmapsto L_1\longmapsto L_{-1}.
\]
This operation fixes \(L_0\) and every selected far field,
including its jets, provided both signs of each far index
were selected.  Indeed, evenness of \(L\) makes the middle
pullback send \(L_j\) to \(L_{-j}\); the reversed path
fixes that selected primary germ, and the last pullback
restores \(L_j\).  The total base pullback is the identity,
so the derivative chain rules cancel.  The operation is
depicted in \Cref{fig:monodromy}.

\begin{figure}[t]
\centering
\begin{tikzpicture}[>=Stealth,node distance=1.0cm,
 state/.style={draw,rounded corners=2pt,minimum width=1.45cm,minimum height=.85cm}]
\node[state] (a) {$L_1$};
\node[state,right=of a] (b) {$M_1$};
\node[state,right=of b] (c) {$M_1$};
\node[state,right=of c] (d) {$L_1$};
\node[state,right=of d] (e) {$L_{-1}$};
\draw[->] (a)--node[above,font=\small] {$A$} (b);
\draw[->] (b)--node[above,font=\small] {$I$} (c);
\draw[->] (c)--node[above,font=\small] {$A^{-1}$} (d);
\draw[->] (d)--node[above,font=\small] {$I$} (e);
\node[font=\small,align=center,below=0.5cm of c]
 {$I:t\mapsto-t$,\qquad $M_1(y,-t)=M_1(y,t)$\\[3pt]
 $L_0$ and the chosen finite collection of far endpoint jets return unchanged.};
\end{tikzpicture}
\caption{The continuation operation in the no-patch argument.
The alternate transverse branch is even at glancing. For each prescribed
finite set of far indices, the path \(A\) is chosen to fix both signs of
those indices. The operation produces another root of the same endpoint
polynomial; it requires no common path for infinitely many indices.}
\label{fig:monodromy}
\end{figure}


Paths and generic basepoints can be chosen anew for each
finite selection.  The identities before and after are
identities of actual primary germs near real glancing,
and hence transport between those basepoints.

Apply \eqref{cont:monodromy-operation} to algebraicity of
\(L_1\) over the jets of \(L_0,L_b\).  It proves algebraicity
of \(L_{-1}\) over the same jets.  Translate this assertion.
From the endpoints of any interval \([B,B+b]\), the line
immediately to its left is now algebraic, as are all
interior lines by \Cref{thm:lens-lines}.  Differentiation
preserves algebraicity over a characteristic-zero
differential field.  Shift the interval one place left
and repeat; transitivity proves algebraicity of every
line to the left over the original endpoint jets.

Choose \(B>K+2\).  There is one fixed nonzero polynomial
over jets of \(L_B,L_{B+b}\) for a generic direction
coordinate of \(L_1\).  For each \(n\ge0\), first shear an
identity by \((y,t)\mapsto(y+2nt,t)\), apply
\eqref{cont:monodromy-operation} with far indices including
\(\pm(B+2n),\pm(B+b+2n)\), and shear back.  It sends
\[
L_{1-2n}\longmapsto L_{-1-2n}
\]
and fixes the original coefficient jets: the shear sends
\(L_{1-2n}\) to \(L_1\) and the coefficient fields to the selected
far indices, while the inverse shear returns those fields exactly.
For \(S_qF(y,t)=F(y+qt,t)\),
\(S_q(\partial_tF)=\partial_t(S_qF)-q\partial_y(S_qF)\);
thus the shear chain rules cancel with those of its inverse.
Inductively this
same polynomial has the direction coordinates of
\[
 L_1,L_{-1},L_{-3},L_{-5},\ldots
\]
as roots in the field of primary germs.
They are distinct: since
\(\Theta(y,t)=\phi(y)+O(t^2)\), the first \(t\)-derivative
of a generic slope coordinate of \(L_j\) is
\(j\phi'(y)\) times a nonzero chart factor.  A nonzero
polynomial over a field cannot have infinitely many
distinct roots.  This contradiction excludes a first
wall at or below \(H_1+\eta\) for generic real irrational
parameters in the neighborhood under consideration.

\paragraph{Transfer to the original real parameter.}
For every nearby irrational \(t_1\), the preceding
construction supplies a parameter neighborhood on which
\[
 h_L(t)\ge H_1+\eta
\]
for real \(t\) outside a countable exceptional set.
The increment \(\eta\) depends only on the initial patch.
A countable subcover covers all nearby irrational real
parameters.  Discarding their exceptional sets and the
countably many rational shifts yields the displayed
bound almost everywhere on an interval about \(t_0\).

Put \(M=H_1+\eta\).  On an upper half disk with diameter
in that interval, approximate the nonnegative lower
semicontinuous boundary values of \(h_L\) from below by
continuous functions supported away from the diameter
endpoints.  Superharmonic comparison and the
almost-everywhere bound give
\[
 h_L(t)\ge M\omega_+(t),
\]
where \(\omega_+\) is the harmonic measure of the diameter.
The lower half disk gives \(h_L\ge M\omega_-\).
Both harmonic measures tend to one near the interior
diameter point \(t_0\).  The small-circle supermean
inequality centered at \(t_0\), followed by shrinking
the circle, gives \(h_L(t_0)\ge M\).  This contradicts
\(h_L(t_0)=H_1\).

\paragraph{The same conclusion for impact patches.}
Suppose now that \(X\) has a joint patch through its
finite first wall at a real parameter.  At real
parameters \eqref{cont:width-equality} gives \(h_X=h_L\).
For nearby irrational shifts, forward and backward
translates
\[
 X(y+(j+1/2)t,t)
\]
of the patch supply enclosing impacts for every edge.
Choose a finite covering by phase boxes, enclosing indices
with at least two steps between them, and a common
narrower vertical band exactly as above.
\Cref{thm:lens-points}, with the analytic coefficient
\(A(t)\), supplies full-band joint polynomial equations
for the edge \(L(y,t)\), attached to the primary branch
on a common bottom strip.

These equations are exactly the input to the first part
of the argument.  Generic slices, root bounds and
monodromy turn a first obstruction into a finite defect.
Its exclusion still uses the line-endpoint identities of
\Cref{thm:lens-lines} and the evenness operation
\eqref{cont:monodromy-operation}; no additional point
identity is needed in that exclusion.  Thus
\(h_X(t)=h_L(t)\ge H_1+\eta\) almost everywhere on a
real interval.  Apply the same superharmonic transfer to
\(h_X\), using \Cref{prop:width-superharmonic}, to
contradict its value at the original parameter.
Real symmetry supplies the lower-wall assertion.
\end{proof}

\subsection{The boundary parametrization has the same natural wall}

The zero slice of the joint impact map is the boundary curve.
We must rule out the possibility that this slice continues
through a wall even though the two-variable family does not.
The argument first propagates glancing trajectories in steep
complex directions and then fills the remaining directions
by a family of quadratic disks.

\begin{corollary}[Natural boundary of the shape]
\label{cor:shape-natural-boundary}
Put \(H=h_X(0)=h_L(0)\).  The parametrization
\(\Gamma(y)=X(y,0)\) is meromorphic on
\[
 D_H=\{y\in\C:|\Im y|<H\}.
\]
If \(H<\infty\), it has no meromorphic continuation patch
through any point of either boundary line of \(D_H\).
Moreover, on its meromorphic domain,
\begin{equation}
 \frac{\det(\Gamma',\Gamma'')^4}
      {(\Gamma'\cdot\Gamma')^3}=c^2,
 \label{cont:intrinsic-normalization}
\end{equation}
where \(c\ne0\) is the normalized curvature constant from
the glancing construction.
\end{corollary}

\begin{proof}
At each strict upper subheight of \(H\), \(X\) has a joint
meromorphic extension near \(t=0\).  In generic affine
coordinates there is no polar component along the entire
zero slice: such a component would meet the initial
collar, where \(X(y,0)=\Gamma(y)\) is regular.
Restrictions to that slice are therefore meromorphic,
with only isolated exceptional points to remove in the
one-variable projective map.  They agree on overlaps by
the identity theorem.  Real symmetry supplies the lower
half, proving the first assertion.

On the original real arc the curvature normalization is
\(r(\phi(y))(\phi'(y))^3=c\).  Since the unit tangent is
\(e_\phi\), it gives
\[
 \Gamma'=\frac{c}{(\phi')^2}e_\phi,\qquad
 \Gamma'\cdot\Gamma'=\frac{c^2}{(\phi')^4},\qquad
 \det(\Gamma',\Gamma'')=\frac{c^2}{(\phi')^3}.
\]
These equations prove \eqref{cont:intrinsic-normalization}
on the real arc, hence meromorphically everywhere on
\(D_H\) and on any continuation patch.

Suppose \(H<\infty\) and a patch exists.  By real symmetry
we may use the upper edge.  Choose a regular point
\[
 s_{\rm edge}=x_*+\ii H
\]
of the patch where \(\Gamma'\cdot\Gamma'\) and
\(\det(\Gamma',\Gamma'')\) are both nonzero.  Such a point
exists after a small move along the edge: all exceptional
sets are isolated, since the corresponding meromorphic
functions are nonzero on the original real arc.
Shrink the patch so that \(\Gamma\), its normal, a local
normal angle \(\phi\), the inverse angle coordinate, and
the curvature radius are holomorphic and nondegenerate.
The normal is obtained from \(\Gamma'\) by a local square
root of its bilinear squared length; the angle then
comes from a local logarithm.  The intrinsic identity
\eqref{cont:intrinsic-normalization} shows that
\(r(\phi)(\phi')^3\) is a nonzero constant on this patch,
equal to \(c\) up to the local sign choices.  The zero
slice \(L(y,0)\) is its tangent, by continuation from below.

\paragraph{A regular starting segment.}
Choose a row slightly below \(H\), sufficiently close
that all short slanted segments with normalized vertical
component of absolute value at least \(1/3\), ending near
\(s_{\rm edge}\), remain inside the regular patch.  Their
starting points lie in a compact horizontal segment.
Choose the exact row height to avoid joint poles and
indeterminacies of \(L,X\) on \(t=0\) over this segment.
This is possible because in affine charts finite along
the segment there is no polar component equal to the
zero slice; joint exceptions intersect it discretely.
Compactness then gives a neighborhood of the segment
on which \(L,X\) are analytic for small complex \(t\).
Evenness gives the line angle expansion
\begin{equation}
 \Theta(y,t)=\phi(y)+O(t^2)
 \label{cont:patch-angle-expansion}
\end{equation}
uniformly there.

The rank-one identity identifies the nearby impact image
with the fixed regular curve arc traced on \(t=0\).
The reflection identity is its ordinary local reflection
law.  Thus the successive line angles, expressed in the
normalized coordinate
\[
 V=\phi^{-1}(\hbox{line angle}),
\]
satisfy the stationary recurrence of the local action
\(D\).  The support and all root branches are chosen
locally; the line angles are close enough that their
mean is the normal angle at the shared impact.

\paragraph{Cubic accuracy of the local recurrence.}
The same local Taylor calculation used for the normalized
action, now using constancy of \(r(\phi)(\phi')^3\), gives
\[
 \partial_1D(V,V+s)=\kappa s^2+O(s^4),
 \qquad \kappa\ne0\ \hbox{constant}.
\]
There is no cubic term.  Antisymmetry of \(D\) turns
stationarity at a node into equality of this expression
for the forward and backward increments.  Its analytic
square root is \(s\sqrt{\kappa+O(s^2)}\).  Choosing the
opposite root from the repeated previous node gives the
nonrepeated successor analytically and uniformly:
\begin{equation}
 V_{j+1}=2V_j-V_{j-1}
                +O\bigl((V_j-V_{j-1})^3\bigr).
 \label{cont:steep-recurrence}
\end{equation}
The offsets of the corresponding lines are also analytic,
being the endpoint momenta of the local generating
function \(\mathcal S\).

\paragraph{Propagation in steep directions.}
Assume \(|\Im t|\ge |t|/3\).  Replace \(t\) by \(-t\)
if necessary, using evenness, so that the step is upward.
For a target \(y\) near the edge choose \(N\) with
\(y-Nt\) in the starting neighborhood.  There is a uniform
bound \(N=O(1/|t|)\).  Initialize the recurrence with
\[
 V_0=y-Nt+O(t^2),\qquad
 V_1=y-(N-1)t+O(t^2),
\]
using \eqref{cont:patch-angle-expansion}.
If the step sizes stay at most \(2|t|\), each recurrence
error is \(O(|t|^3)\).  Summing the errors once gives
\[
 V_j-V_{j-1}=t+O(t^2),
\]
and summing a second time gives
\[
 V_j-\bigl(y-(N-j)t\bigr)=O(|t|),
 \qquad 0\le j\le N.
\]
These estimates close the assumed step bound for small
\(|t|\).  Choose the geometric paths with fixed slack in
the regular patch before shrinking \(|t|\); then every
iterate remains in that patch.

For locally constant \(N\), this construction is
holomorphic in \(y,t\).  Choices of different \(N\)'s
agree: take their starting points in a thin slab of
known analytic data, where the exact recurrence relates
all intervening indices, then apply uniqueness of the
nonrepeated successor.  It follows that the resulting
germs glue near the target and along a vertical corridor
down to the starting slab, uniformly for all steep
directions.  They are attached to the primary data on
a fixed lower part of that corridor.  At any chosen
strict subheight below \(H\), make \(|t|\) smaller if
needed so that uniqueness identifies them with the
already known primary branch.  The line construction
also gives \(X\) by the neighboring intersections.
All choices have been made in order: first the regular
patch and geometric slack, then the vertical travel
distance, and finally the allowed size of \(|t|\).

\paragraph{Quadratic disks fill the shallow directions.}
Choose \(\eta>0\) small compared with both the corridor
size and the depth of the starting row.  Next choose
a sufficiently small parameter radius \(r_0>0\).  Consider
the analytic disks
\begin{equation}
 t=r_0\zeta,\qquad
 y=x_*+\ii\sigma-\ii\eta\zeta^2+w,\qquad
 H-2\eta\le\sigma\le H,
 \label{cont:quadratic-disks}
\end{equation}
where \(w\) is a small complex transverse parameter.
At \(|\zeta|=1\), write \(\zeta=e^{\ii\theta}\).
On the near-real arcs \(|\sin\theta|<1/2\),
\[
 \Im y\big|_{w=0}
   =\sigma-\eta\cos(2\theta)
   \le H-\eta/2.
\]
These arcs therefore have joint primary data for small
\(r_0\).  On the remaining arcs,
\(|\Im t|/|t|\ge1/2\), so the steep continuation just
constructed applies with a strict margin over \(1/3\).
The continued germs glue to the primary ones on the
overlaps by their attachment below \(H\).
Compactness supplies a common annular collar in \(\zeta\)
and a common small disk in \(w\) on which the pullbacks
are meromorphic, throughout the indicated interval of
\(\sigma\).

At \(\sigma=H-2\eta\), the full disks have height at most
\(H-\eta\) for \(w=0\).  Hence they, and a small
neighborhood of them, already lie in the primary joint
domain when \(r_0\) is sufficiently small.
Increase \(\sigma\) in small steps.  The displacement of
one center is absorbed by the transverse \(w\)-disk, so
the full disks from the previous step give a nonempty
open seed at the next step.  Apply
\Cref{lem:meromorphic-hartogs,cont:disc-sweep} at each
step.  Uniqueness on the collars and on the seed
overlaps attaches the resulting extensions to the lower
branch throughout the ascent.

At \(\sigma=H\), the disk center \(\zeta=0\) gives a
joint meromorphic patch through
\((y,t)=(x_*+\ii H,0)\); varying \(w\) and \(\zeta\)
are local coordinates there because \(r_0\ne0\).
This contradicts \Cref{thm:no-patch}, applied to \(L\)
at the real parameter \(0\).  The supposed shape patch
cannot exist.  Real symmetry excludes a lower-edge
patch as well.
\end{proof}
